% cv_extended.tex — Extended Computer Vision analysis
% Appendix I

\section{Computer Vision: Extended Analysis}\label{app:cv-extended}

The main body of this report (Section~\ref{sec:cv}) summarises the detection pipeline at a level sufficient to understand the system architecture.  This appendix provides the full technical depth: stage-by-stage timing breakdowns measured on the Raspberry Pi~5, comprehensive benchmark results across all available models, a mission-level detection probability analysis that links frame rate to flyover reliability, altitude-dependent performance modelling validated against DJI flight-video replay, and the multi-backend abstraction that enables the same codebase to run transparently on laptop, Pi, and optimised ARM deployments.  These details are essential for reproducing the system, tuning configuration parameters for different mission profiles, and understanding the margins available for future optimisation.

% ═══════════════════════════════════════════════════════════════
\subsection{Detection Pipeline: Step-by-Step Timing}\label{app:cv:pipeline}

Every iteration of the main loop executes the detection pipeline in \texttt{vision.py} through six sequential stages.  Figure~\ref{fig:cv_pipeline} illustrates the data flow; Table~\ref{tab:pipeline_timing} provides measured timing for each stage on the Raspberry Pi~5.

\begin{table}[ht]
  \centering
  \caption{Per-stage timing breakdown for the detection pipeline on Raspberry Pi~5 (mean of 50 runs, YOLOv8n FP32, XNNPACK 4-thread).}\label{tab:pipeline_timing}
  \begin{tabular}{@{}r l c l@{}}
    \toprule
    \textbf{Stage} & \textbf{Operation} & \textbf{Time (ms)} & \textbf{Notes} \\
    \midrule
    1 & Camera capture (picamera2) & $\sim$2 & CSI-2 DMA, $1456\!\times\!1088$ BGR \\
    2 & Lens undistortion (\texttt{cv2.remap}) & 1.5 & Precomputed maps, RMS\,=\,0.399\,px \\
    3 & Preprocessing (cvtColor + resize + norm.) & $\sim$3.5 & BGR$\to$RGB, bilinear to $640\!\times\!640$, $\div 255$ \\
    4 & TFLite inference (XNNPACK) & $\sim$196 & 4 threads on Cortex-A76 cores \\
    5 & Postprocessing (confidence filter + NMS) & $\sim$2 & Greedy single-target selection \\
    6 & Bounding box overlay (\texttt{cv2.rectangle}) & $<$1 & Drawn onto original frame \\
    \midrule
    & \textbf{Total (with undistortion)} & \textbf{$\sim$208} & \textbf{4.8\,FPS effective} \\
    & \textbf{Total (without undistortion)} & \textbf{$\sim$206.5} & \textbf{Undistortion cost: +1.5\,ms} \\
    \bottomrule
  \end{tabular}
\end{table}

\subsubsection{Stage 1: Camera Capture}

The IMX296 global shutter sensor captures frames at $1456 \times 1088$ pixels via CSI-2 using \texttt{picamera2}.  A critical hardware-specific detail: despite being configured as \texttt{RGB888}, the IMX296 outputs data in BGR channel order---identical to OpenCV's native format.  This was confirmed empirically by testing all six channel permutations against known colour targets.  Consequently, no \texttt{cv2.cvtColor(frame, cv2.COLOR\_RGB2BGR)} conversion is applied, saving approximately \SI{1.2}{\milli\second} per frame and avoiding a subtle but mission-critical colour inversion that would degrade detection accuracy.

At startup, 10 frames are discarded to allow the auto-white-balance (AWB) and auto-exposure algorithms to converge.  The camera can be configured for preset AWB modes (daylight, cloudy, indoor) or fixed manual colour gains via \texttt{config.py}, depending on lighting conditions.

\subsubsection{Stage 2: Lens Undistortion}

If a lens calibration file (\texttt{calibration\_data.npz}) is present, precomputed remap matrices are loaded at initialisation via \texttt{cv2.initUndistortRectifyMap()}.  Per-frame undistortion uses \texttt{cv2.remap()} with bilinear interpolation, costing \SI{1.5}{\milli\second}---less than 1\% of the total pipeline latency.  The calibration was performed using a 14$\times$9 checkerboard (13$\times$8 inner corners, \SI{28}{\milli\metre} square size), achieving an RMS reprojection error of 0.399 pixels.  This corrects barrel distortion from the wide-angle lens, which otherwise displaces peripheral detections by up to 8 pixels, corresponding to approximately \SI{0.7}{\metre} GPS error at \SI{35}{\metre} altitude.

\subsubsection{Stage 3: Preprocessing}

The undistorted $1456 \times 1088$ BGR frame is converted to RGB (\texttt{cv2.cvtColor}), resized to $640 \times 640$ via bilinear interpolation (\texttt{cv2.resize}), cast to \texttt{float32}, and normalised to $[0, 1]$ by dividing by 255.  The result is expanded to a batch tensor of shape $[1, 640, 640, 3]$.

Direct resize (without letterboxing) was chosen over the alternative of maintaining aspect ratio with padding for three reasons: (1)~the aspect ratio of the IMX296 sensor ($4{:}3$) is close to square, so the distortion introduced by non-uniform scaling is modest (${\sim}7\%$ vertical stretch); (2)~letterboxing wastes ${\sim}14\%$ of the $640 \times 640$ input area on grey padding, reducing the effective resolution available for small-target detection; (3)~the coordinate rescaling from model output back to original frame dimensions is a simple linear map, avoiding the padding-offset arithmetic that introduces rounding errors.

\subsubsection{Stage 4: TFLite Inference}

The preprocessed tensor is passed to the TFLite interpreter, which invokes the XNNPACK delegate with four threads (one per Cortex-A76 core).  This stage dominates the pipeline at \SI{196}{\milli\second} (94\% of total latency).

The YOLOv8n architecture produces an output tensor of shape $[1, 5, 8400]$, where:
\begin{itemize}
    \item \textbf{8400} is the total number of anchor-free detection candidates, arising from three feature pyramid network (FPN) scales: $80 \times 80 = 6400$ candidates at the finest scale (detecting small objects), $40 \times 40 = 1600$ at medium scale, and $20 \times 20 = 400$ at the coarsest scale (detecting large objects).
    \item \textbf{5} encodes $[x_\mathrm{centre}, y_\mathrm{centre}, w, h, \mathrm{class\_conf}]$ for a single-class model.  For multi-class models (e.g.\ the 80-class COCO backup model), this dimension extends to $4 + C$ where $C$ is the number of classes.
\end{itemize}

The output tensor is transposed from $[1, 5, 8400]$ to $[8400, 5]$ for row-wise iteration during postprocessing.

\subsubsection{Stage 5: Postprocessing}

The postprocessing stage applies a confidence threshold of 0.2 (configurable via \texttt{config.py}) and selects the single highest-confidence detection.  This threshold was deliberately set well below the typical YOLO default of 0.5 for a critical operational reason: in a search-and-rescue context, a \emph{missed detection is irrecoverable}---the drone may never revisit that ground patch---while a \emph{false positive} is filtered by the human operator at the VERIFY stage and costs only a brief investigation diversion.  The asymmetric cost of errors therefore favours high recall (low threshold) over high precision (high threshold).

For the custom single-class model, no non-maximum suppression (NMS) is required because the system uses a greedy best-detection strategy: a single target per frame is sufficient for the state machine's decision logic.  When the Ultralytics backend is used (laptop development), NMS is handled internally by the framework.

Coordinates are rescaled from the $640 \times 640$ model space to the original $1456 \times 1088$ frame dimensions via simple multiplication: $x_\mathrm{frame} = x_\mathrm{model} \times (1456 / 640)$.

\subsubsection{Stage 6: Visualisation}

A green bounding box with a confidence label (e.g.\ ``\texttt{AI 0.97 [dummy]}'') is drawn onto the original frame using \texttt{cv2.rectangle} and \texttt{cv2.putText}.  This annotated frame is simultaneously served to the ground station MJPEG stream and (in non-headless mode) displayed locally via \texttt{cv2.imshow}.

\begin{figure}[H]
\centering
\includegraphics[width=\linewidth]{vision_pipeline}
\caption{End-to-end vision pipeline: from raw CSI-2 capture through lens undistortion, resize, TFLite inference, confidence filtering, and GPS coordinate projection. Stage timings are measured on the Raspberry Pi~5; total latency is \SI{208}{\milli\second} per frame (Table~\ref{tab:pipeline_timing}).}
\label{fig:vision_pipeline}
\end{figure}

% ═══════════════════════════════════════════════════════════════
\subsection{Comprehensive Benchmark Results}\label{app:cv:benchmarks}

Table~\ref{tab:cv_benchmarks} presents the full benchmark results from the Pi~5, measured using the standardised protocol described in Section~\ref{sec:cv:inference:method} (5 warm-up frames discarded, 50 timed passes, \texttt{time.perf\_counter()}).

\begin{table}[ht]
  \centering
  \caption{Inference benchmarks on Raspberry Pi~5 (Cortex-A76 $\times$4 at \SI{2.4}{\giga\hertz}, XNNPACK 4-thread, active cooler, no desktop environment).}\label{tab:cv_benchmarks}
  \begin{tabular}{@{}l c c c c c c@{}}
    \toprule
    \textbf{Model} & \textbf{Size} & \textbf{Undistort} & \textbf{Mean (ms)} & \textbf{FPS} & \textbf{Det. rate} & \textbf{Mean conf.} \\
    \midrule
    best.tflite (YOLOv8n custom)     & 3.2\,MB  & No  & 206.5 & 4.8 & 50/50 (100\%) & 0.966 \\
    best.tflite (YOLOv8n custom)     & 3.2\,MB  & Yes & 208.0 & 4.8 & 50/50 (100\%) & 0.958 \\
    sar\_v2\_1088 (retrained v2)     & 11.7\,MB & No  & 206.8 & 4.8 & 50/50 (100\%) & 0.971 \\
    human.tflite (COCO 80-class)     & 13.0\,MB & No  & 207.2 & 4.8 & 43/50 (86\%) & 0.724 \\
    \bottomrule
  \end{tabular}
\end{table}

Several observations warrant discussion:

\paragraph{Inference time is model-size-independent.}
All three models produce nearly identical inference latencies (206--208\,ms) despite a $4\times$ size difference (3.2\,MB vs 13.0\,MB).  This confirms that YOLOv8n inference is compute-bound (limited by NEON multiply-accumulate throughput) rather than memory-bound (limited by weight loading from DRAM).  The model architecture---not the weight count per class---determines latency, because all three models share the identical YOLOv8n backbone; only the final classification head differs.

\paragraph{Undistortion overhead is negligible.}
Enabling lens undistortion adds only \SI{1.5}{\milli\second} ($<$1\% overhead), confirming that precomputed remap matrices are an efficient correction strategy.  The slight confidence decrease with undistortion (0.966 $\to$ 0.958) is within measurement noise and does not affect detection rate.

\paragraph{Custom model vastly outperforms generic COCO.}
The custom SAR model achieves 100\% detection rate with 0.966 mean confidence, compared to 86\% and 0.724 for the generic COCO person detector.  This validates the domain-specific retraining strategy (Section~\ref{sec:training}): a model trained on aerial views of a mannequin against grassy backgrounds significantly outperforms a general-purpose detector trained on ground-level person images.  The COCO model is retained as a backup for scenarios where the target is a standing human rather than a prone mannequin.

\paragraph{Timing stability.}
The standard deviation of per-frame latency was \SI{3.2}{\milli\second} ($\sigma / \mu = 1.5\%$), indicating highly reproducible timing.  No thermal throttling was observed during the 50-frame test window with the official Pi~5 active cooler.

% ═══════════════════════════════════════════════════════════════
\subsection{Mission Runtime Analysis}\label{app:cv:runtime}

The practical question for mission planning is not ``what is the per-frame detection rate?'' but rather ``what is the probability of detecting a target during a single flyover of the search swath?''  This subsection models that probability using the benchmark data and mission geometry.

\subsubsection{Ground Footprint Geometry}

At search altitude $h$, the camera's ground footprint is determined by the horizontal and vertical fields of view:
\begin{align}
    W_g &= 2h \tan\!\left(\frac{\theta_H}{2}\right) \label{eq:footprint_w} \\
    L_g &= 2h \tan\!\left(\frac{\theta_V}{2}\right) \label{eq:footprint_l}
\end{align}
where $\theta_H = 2\arctan(w_s / 2f) \approx 49.3^{\circ}$ is the horizontal FOV (sensor width $w_s = \SI{5.02}{\milli\metre}$, focal length $f = \SI{5.46}{\milli\metre}$) and $\theta_V \approx 37.9^{\circ}$ is the vertical FOV.

Table~\ref{tab:footprint} gives the ground footprint at key operational altitudes.

\begin{table}[ht]
  \centering
  \caption{Camera ground footprint and target size at operational altitudes (IMX296, $f = \SI{5.46}{\milli\metre}$, $1456\!\times\!1088$ sensor).  Sensor px = height in native frame; Model px = after resize to $640\!\times\!640$ (scale $= 640/1456 = 0.4396$).}\label{tab:footprint}
  \begin{tabular}{@{}c c c c c c@{}}
    \toprule
    \textbf{Altitude (m)} & $W_g$ \textbf{(m)} & $L_g$ \textbf{(m)} & \textbf{Sensor (px)} & \textbf{Model (px)} & \textbf{GSD (mm/px)} \\
    \midrule
    10  & 9.2  & 6.9  & 285 & 125 & 6.3 \\
    15  & 13.8 & 10.3 & 190 & 84  & 9.5 \\
    20  & 18.4 & 13.8 & 142 & 63  & 12.7 \\
    25  & 23.0 & 17.2 & 114 & 50  & 15.8 \\
    30  & 27.6 & 20.7 & 95  & 42  & 19.0 \\
    35  & 32.2 & 24.1 & 81  & 36  & 22.1 \\
    40  & 36.8 & 27.6 & 71  & 31  & 25.3 \\
    50  & 46.0 & 34.5 & 57  & 25  & 31.6 \\
    \bottomrule
  \end{tabular}
\end{table}

The ``Sensor px'' column shows the dummy's height in the native $1456 \times 1088$ frame ($= 1.8\,\text{m} / \mathrm{GSD}$, where $\mathrm{GSD} = L_g / 1088$).  The ``Model px'' column applies the resize scale factor $640/1456 = 0.4396$ to convert to the $640 \times 640$ inference input.

\subsubsection{Frames per Flyover}

At search speed $v$ and frame rate $r$, the along-track distance between consecutive frames is $d_f = v / r$.  The number of frames in which a stationary target remains within the field of view during a single pass is:
\begin{equation}
    N_\mathrm{frames} = \frac{L_g}{d_f} = \frac{L_g \cdot r}{v}
    \label{eq:frames_per_flyover}
\end{equation}

Table~\ref{tab:frames_flyover} evaluates this for the operational speed and altitude envelope.

\begin{table}[ht]
  \centering
  \caption{Frames per flyover and detection probability for a target within the search swath.  $p_d$ = single-frame detection probability; $P_\mathrm{miss}$ = $(1 - p_d)^N$; $P_\mathrm{detect}$ = $1 - P_\mathrm{miss}$.}\label{tab:frames_flyover}
  \begin{tabular}{@{}c c c c c c@{}}
    \toprule
    \textbf{Alt. (m)} & \textbf{Speed (m/s)} & $N_\mathrm{frames}$ & $p_d$ & $P_\mathrm{detect}$ & \textbf{Note} \\
    \midrule
    35 &  6 & 19.3 & 0.90 & $>$0.9999 & Slow search \\
    35 &  8 & 14.5 & 0.90 & $>$0.9999 & Normal search \\
    35 & 10 & 11.6 & 0.90 & 0.9997 & Default \texttt{SEARCH\_SPEED} \\
    35 & 15 &  7.7 & 0.90 & 0.9974 & Transit speed \\
    50 & 10 & 16.6 & 0.70 & 0.9998 & High altitude \\
    50 & 15 & 11.0 & 0.70 & 0.9983 & High alt + fast \\
    20 & 10 &  6.6 & 0.98 & $>$0.9999 & Low altitude pass \\
    15 &  5 &  9.9 & 0.99 & $>$0.9999 & Centering descent \\
    \bottomrule
  \end{tabular}
\end{table}

\paragraph{Key finding.}
Even in the worst operational case (50\,m altitude at 15\,m/s transit speed, with a conservative 70\% single-frame detection probability), the cumulative detection probability across the 11-frame flyover window exceeds 99.8\%.  At the nominal search configuration (35\,m, 8\,m/s from the altitude-dependent schedule), the system achieves effectively certain detection ($P > 99.99\%$) with 14.5 frames per target.

The lawnmower pattern further improves robustness: adjacent passes overlap by 20\% (configurable via \texttt{planning.py}), so targets near the swath edge receive a second observation window.  The probability of missing a target across two overlapping passes at nominal conditions is $(1 - 0.9997)^2 \approx 10^{-7}$.

\subsubsection{Frame Overlap Analysis}

At 4.8\,FPS and the nominal 8\,m/s search speed, consecutive frames are separated by $8/4.8 = \SI{1.67}{\metre}$ along-track.  The along-track footprint at \SI{35}{\metre} is \SI{24.1}{\metre}, giving a frame-to-frame overlap of:
\begin{equation}
    \mathrm{Overlap} = 1 - \frac{d_f}{L_g} = 1 - \frac{1.67}{24.1} = 93.1\%
\end{equation}

This high overlap means that a target appears in many consecutive frames, providing the state machine with multiple opportunities to trigger the CENTERING sequence.  The \texttt{-{}-smart-detect} mode exploits this by requiring $k = 3$ consecutive positive detections before committing to investigation, effectively implementing a temporal majority-vote filter that suppresses transient false positives at a cost of only $3 \times 208 = \SI{0.6}{\second}$ delay (\SI{5.0}{\metre} of travel at 8\,m/s).

% ═══════════════════════════════════════════════════════════════
\subsection{Detection Performance vs Altitude}\label{app:cv:altitude}

The critical operational question is: ``at what altitude does detection become unreliable?''  This is modelled using the thin-lens equation to predict the target's apparent size in pixels, combined with empirically observed confidence behaviour from video replay analysis. Figure~\ref{fig:gsd_altitude} visualises the relationship between altitude, ground sampling distance, and target pixel size.

\begin{figure}[H]
\centering
\includegraphics[width=0.8\linewidth]{gsd_altitude}
\caption{Ground sampling distance (GSD) and target pixel size as a function of altitude for the IMX296 camera ($f = \SI{5.46}{\milli\metre}$, $w_s = \SI{5.02}{\milli\metre}$, $1456 \times 1088$\,px). The horizontal dashed line marks the 20-pixel detection floor in model-input space; above \SI{63}{\metre} the \SI{1.8}{\metre} mannequin falls below this threshold.}
\label{fig:gsd_altitude}
\end{figure}

\subsubsection{Pixel Size Model}

The target of real height $H_t = \SI{1.8}{\metre}$ first projects onto the sensor with $\mathrm{px}_{\mathrm{sensor}} = H_t \cdot f_{\mathrm{px}} / h$ pixels, where $f_{\mathrm{px}} = f / w_s \times W_{\mathrm{px}} = 5.46/5.02 \times 1456 = 1584$\,px.
The $1456 \times 1088$ frame is then resized to $640 \times 640$; applying the uniform scale factor $640/1456 = 0.4396$ gives the model-input height:
\begin{equation}
    p_t = \frac{H_t \cdot f_{\mathrm{px}}}{h} \times \frac{640}{1456}
    \label{eq:target_pixels}
\end{equation}

\subsubsection{Altitude Performance Table}

Table~\ref{tab:altitude_performance} combines the pixel-size model with confidence data extracted from DJI video replay analysis (Section~\ref{sec:localisation}) and bench testing.

\begin{table}[ht]
  \centering
  \caption{Detection performance vs altitude for a \SI{1.8}{\metre} mannequin target.  Pixel sizes are shown for both the native sensor frame and the $640\!\times\!640$ model input (scale $= 640/1456 = 0.4396$).  Confidence and detection rate are empirically observed or interpolated from video replay data.}\label{tab:altitude_performance}
  \begin{tabular}{@{}c c c c c p{4.5cm}@{}}
    \toprule
    \textbf{Alt.\ (m)} & \textbf{Sensor (px)} & \textbf{Model (px)} & \textbf{Conf.} & \textbf{Det.\ rate} & \textbf{Assessment} \\
    \midrule
    10  & 285 & 125 & 0.98 & 100\% & Excellent --- target fills $>$19\% of model height \\
    15  & 190 & 84  & 0.97 & 100\% & Excellent --- verify altitude \\
    20  & 142 & 63  & 0.96 & 100\% & Very good --- still large and clear \\
    25  & 114 & 50  & 0.94 & 98\%  & Good --- typical low search altitude \\
    30  & 95  & 42  & 0.92 & 96\%  & Good --- reliable detection \\
    35  & 81  & 36  & 0.90 & 95\%  & Good --- nominal search altitude \\
    40  & 71  & 31  & 0.85 & 90\%  & Marginal --- approaching limits \\
    50  & 57  & 25  & 0.70 & 75\%  & Degraded --- single-pass detection unreliable \\
    60  & 48  & 21  & 0.45 & 50\%  & Poor --- degraded confidence \\
    70  & 41  & 18  & 0.30 & 30\%  & Unreliable --- approaching confidence threshold \\
    80  & 36  & 16  & $<$0.20 & $<$15\% & Not viable --- below confidence threshold \\
    \bottomrule
  \end{tabular}
\end{table}

\begin{figure}[ht]
  \centering
  \includegraphics[width=0.85\linewidth]{altitude_detection_table}
  \caption{Detection performance as a function of altitude for the YOLOv8n custom model against a \SI{1.8}{\metre} mannequin target. The table combines predicted pixel sizes from the thin-lens model with empirically observed confidence and detection rates from video replay analysis.}
  \label{fig:altitude-detection-table}
\end{figure}

\paragraph{Operational envelope.}
The data confirms that \SI{35}{\metre} (the configured \texttt{TARGET\_ALT}) sits comfortably within the reliable detection zone (95\% per-frame, $>$99.97\% per-flyover).  The altitude ceiling for reliable single-pass detection is approximately \SI{40}{\metre}; above this, multi-pass coverage from the lawnmower pattern provides the necessary redundancy.  The theoretical floor for detection (${\sim}\SI{70}{\metre}$) corresponds to the target subtending approximately 18 model-input pixels (41 sensor pixels)---near the YOLOv8n smallest feature scale ($80 \times 80$ grid, where each cell covers an $8 \times 8$ region of the input).

\paragraph{Confidence versus pixel size relationship.}
Figure~\ref{fig:conf_vs_alt_ext} shows the empirically measured confidence as a function of altitude.  The relationship follows an approximately sigmoidal curve in log-altitude space, with a sharp transition between reliable detection ($>$0.9) and unreliable detection ($<$0.5) occurring over a narrow altitude band (\SIrange{40}{60}{\metre}).  This sharp transition is characteristic of anchor-free detectors operating near their minimum detectable object size.

\begin{figure}[ht]
\centering
\includegraphics[width=0.75\textwidth]{conf_vs_alt}
\caption{Detection confidence vs altitude for the YOLOv8n custom model.  Error bars represent the interquartile range from video replay analysis.  The dashed line marks the 0.2 confidence threshold.}
\label{fig:conf_vs_alt_ext}
\end{figure}

\paragraph{Combined detection envelope.}
Figure~\ref{fig:detection_envelope} overlays the altitude-dependent confidence curve with the motion blur degradation model, producing a combined detection envelope that accounts for both target size reduction and image quality loss. The intersection of these two degradation factors defines the practical operational ceiling more accurately than either factor alone.

\begin{figure}[ht]
\centering
\includegraphics[width=0.85\textwidth]{detection_envelope}
\caption{Combined detection envelope showing confidence degradation from both altitude (target pixel size reduction) and motion blur at \SI{10}{\metre\per\second}. The shaded region indicates the reliable detection zone where both factors remain within acceptable limits. The nominal search altitude of \SI{35}{\metre} sits comfortably within this envelope.}
\label{fig:detection_envelope}
\end{figure}

% ═══════════════════════════════════════════════════════════════
\subsection{Model Training Results}\label{app:cv:training}

The retrained v2 model (\texttt{sar\_v2\_1088}) was trained on Google Colab (NVIDIA T4 GPU) for 150 epochs using the mixed dataset described in Section~\ref{sec:training}.  Table~\ref{tab:training_results_extended} summarises the final evaluation metrics.

\begin{table}[ht]
  \centering
  \caption{Training evaluation metrics for the YOLOv8n SAR detector (v2, trained at $1088 \times 1088$, evaluated on validation split).}\label{tab:training_results_extended}
  \begin{tabular}{@{}l c l@{}}
    \toprule
    \textbf{Metric} & \textbf{Value} & \textbf{Interpretation} \\
    \midrule
    mAP$_{50}$         & 0.995 & Near-perfect at IoU $\geq 0.50$ \\
    mAP$_{50{:}95}$    & 0.828 & Strong across IoU thresholds \\
    Precision          & 0.994 & Very low false positive rate \\
    Recall             & 0.989 & Near-zero missed detections \\
    F1 score           & 0.992 & Balanced precision-recall \\
    \midrule
    Training time      & 47\,min & 150 epochs on T4 GPU \\
    Best epoch         & 127 & Selected by early stopping (patience=30) \\
    Final model size   & 11.7\,MB (TFLite FP32) & Input $[1,640,640,3]$, output $[1,5,8400]$ \\
    \bottomrule
  \end{tabular}
\end{table}

\paragraph{Confusion matrix analysis.}
The single-class confusion matrix on the validation set shows 99.5\% true positive rate and 0.5\% false negative rate.  No false positives were recorded on the 50-image negative set (background frames without mannequin), confirming that the negative-mining strategy successfully suppresses hallucinated detections on shadows, paths, and field markings that plagued the v1 model.

\paragraph{Blur-robustness through data augmentation.}
A critical design decision in the training pipeline was the inclusion of augmentations that simulate real flight conditions.  The dataset generator (\texttt{generate\_dataset\_v2.py}) applies four offline augmentations: random \emph{brightness} jitter ($\alpha \in [0.7, 1.3]$), \emph{contrast} offset ($\beta \in [-30, 30]$), \emph{Gaussian blur} ($3\!\times\!3$ or $5\!\times\!5$ kernel, 30\% probability), and random \emph{rotation} ($0$--$360^{\circ}$).  These are complemented by the YOLO trainer's online augmentations---mosaic, scale, mixup, copy-paste, and horizontal flip (Table~\ref{tab:augmentation})---which further diversify the training distribution.  Together, these augmentations expose the model to degraded images representative of the conditions analysed in Section~\ref{sec:vp:blur}, ensuring the detector maintains high confidence even when the real camera produces mildly blurred frames due to wind gusts or rapid manoeuvres---a robustness that would not be achieved by training exclusively on sharp synthetic composites.

\paragraph{Comparison with v1.}
The v1 model (trained on 200 synthetic images at $640 \times 640$ from a single satellite background) achieved mAP$_{50}$ ${\sim}0.95$ and suffered frequent false positives on high-contrast terrain features.  The v2 improvements---real-flight backgrounds, native-resolution training, altitude-correct scaling, and explicit negatives---closed the domain gap and improved mAP$_{50}$ by approximately 4.5 percentage points while eliminating the false positive problem.

\paragraph{Export and deployment.}
The trained \texttt{best.pt} weights were exported to TFLite (float32) via \texttt{yolo export format=tflite}, producing the same $[1, 640, 640, 3]$ input and $[1, 5, 8400]$ output shape as the v1 model.  This ensures drop-in compatibility: deploying the v2 model on the Pi requires only copying the \texttt{.tflite} file---no code changes, no configuration changes, no recompilation.

% ═══════════════════════════════════════════════════════════════
\subsection{Dual Backend Architecture}\label{app:cv:backends}

The \texttt{VisionSystem} class in \texttt{vision.py} implements a three-tier backend selection strategy that automatically adapts to the available runtime environment.  This design enables the same codebase to run on a development laptop with full Ultralytics/PyTorch, a Raspberry Pi with lightweight TFLite, or an optimised ARM deployment with NCNN---without any code changes or conditional imports in the calling modules.

\subsubsection{Backend Selection Logic}

At import time, \texttt{vision.py} probes for available backends in priority order:

\begin{enumerate}
    \item \textbf{NCNN} (if \texttt{-{}-backend ncnn} flag is set and the \texttt{ncnn} package is installed).  This backend provides the fastest ARM inference (${\sim}$15\,FPS projected) through hand-tuned NEON assembly kernels, but requires manual compilation on Python~3.13/aarch64 and was therefore not used during field testing.
    \item \textbf{Ultralytics YOLO} (if the \texttt{ultralytics} package is importable).  This is the default on the development laptop, providing the richest feature set (built-in NMS, class name mapping, training integration) at the cost of a large dependency footprint ($>$\SI{800}{\mega\byte}).
    \item \textbf{TFLite} (if \texttt{tflite-runtime}, \texttt{ai-edge-litert}, or \texttt{tensorflow.lite} is available).  This is the production backend on the Pi, adding only ${\sim}$\SI{5}{\mega\byte} to the deployment.  On Python~3.13, the legacy \texttt{tflite-runtime} package is unavailable; the system falls back to \texttt{ai-edge-litert} (Google's successor package).
\end{enumerate}

All three backends consume the same \texttt{.tflite} model file (or equivalent format) and produce output through the identical \texttt{detect\_in\_image(frame)} $\to$ \texttt{(found, cx, cy, confidence)} interface.  The calling code---state machine, GPS estimation, ground station---is entirely backend-agnostic.

\subsubsection{Backend Abstraction Benefits}

This architecture provides three concrete benefits:

\begin{itemize}
    \item \textbf{Development-production parity.}  Detections observed during laptop simulation are reproduced identically on the Pi (modulo floating-point precision differences between CPU architectures), because the model weights and preprocessing pipeline are shared.
    \item \textbf{Graceful degradation.}  If the preferred backend fails to load (e.g.\ missing NCNN compilation), the system automatically falls back to the next available option rather than crashing.  A missing model file produces a clear warning (``Detection DISABLED --- mission will fly but never detect targets'') rather than a silent failure.
    \item \textbf{A/B testing.}  Different backends can be compared on identical input frames using the \texttt{-{}-backend} flag, enabling systematic evaluation of accuracy and latency trade-offs during development.
\end{itemize}

% ═══════════════════════════════════════════════════════════════
\subsection{Model Swapping and Field Deployment}\label{app:cv:swapping}

A key operational requirement is the ability to swap detection models in the field without modifying code or restarting the mission planning workflow.  Two mechanisms are provided:

\paragraph{File-copy swap.}
All scripts read their model from a single path: \texttt{best.tflite} in the project root.  Swapping models is a single copy command:
\begin{lstlisting}[language=bash,numbers=none]
cp cv_models/sar_v2_1088/best.tflite best.tflite
\end{lstlisting}
Three model variants are maintained in the \texttt{cv\_models/} directory (Table~\ref{tab:model_variants}), each with TFLite, PyTorch (\texttt{.pt}), and NCNN exports:

\begin{table}[ht]
  \centering
  \caption{Available model variants for field deployment.}\label{tab:model_variants}
  \begin{tabular}{@{}l c c l@{}}
    \toprule
    \textbf{Model} & \textbf{Size} & \textbf{Classes} & \textbf{Use case} \\
    \midrule
    \texttt{sar\_v2\_1088} & 11.7\,MB & 1 (dummy) & Primary: retrained on real + synthetic data \\
    \texttt{sar\_640}       & 3.2\,MB  & 1 (dummy) & Baseline: original training at $640\!\times\!640$ \\
    \texttt{human.tflite}   & 13.0\,MB & 80 (COCO) & Backup: generic person detector \\
    \bottomrule
  \end{tabular}
\end{table}

\paragraph{CLI flag swap.}
The \texttt{-{}-model} flag overrides the default model path at launch:
\begin{lstlisting}[language=bash,numbers=none]
python main.py --model models/human.tflite
\end{lstlisting}
This enables rapid A/B testing during flight day without touching the filesystem.

\paragraph{Zero-code-change guarantee.}
All TFLite models share the same input tensor shape $[1, 640, 640, 3]$ and output shape $[1, 5, 8400]$ (or $[1, 4+C, 8400]$ for multi-class).  The \texttt{\_pick\_best\_detection()} method in \texttt{vision.py} handles both single-class and multi-class outputs transparently by scanning the class dimension and selecting the maximum confidence.  For multi-class models, the \texttt{COCO\_NAMES} dictionary maps class indices to human-readable labels (e.g.\ ``person'', ``car'') for the bounding box overlay.

% ═══════════════════════════════════════════════════════════════
\subsection{Detection Speed vs Drone Speed}\label{app:cv:speed}

At higher search speeds, two effects degrade detection quality: (1)~fewer frames per flyover, reducing the cumulative detection probability; and (2)~increased motion blur, reducing per-frame confidence.

\paragraph{Frames per flyover.}
Equation~\ref{eq:frames_per_flyover} shows that $N_\mathrm{frames} \propto 1/v$.  At the maximum transit speed of \SI{15}{\metre\per\second}, the flyover window at \SI{35}{\metre} altitude shrinks to 7.7 frames---still sufficient for $P_\mathrm{detect} > 99.7\%$ (Table~\ref{tab:frames_flyover}).

\paragraph{Motion blur.}
At \SI{10}{\metre\per\second} and 4.8\,FPS, each frame integrates over $10 / 4.8 = \SI{2.08}{\metre}$ of ground motion.  With the IMX296's global shutter and a typical exposure time of ${\sim}\SI{2}{\milli\second}$ (auto-exposure at outdoor illumination levels), the motion smear is $0.002 \times 10 = \SI{0.02}{\metre}$---less than 1 pixel at \SI{35}{\metre} altitude (GSD = \SI{22.1}{\milli\metre}/px).  The global shutter eliminates the rolling-shutter ``jello'' artefact that would otherwise cause geometric distortion of the target at speed, making the IMX296 an ideal choice for aerial detection from a moving platform.

Figure~\ref{fig:blur_vs_altitude} quantifies pixel motion blur as a function of altitude for several speed and exposure combinations, confirming that the IMX296's short exposure times keep blur well below one pixel across the entire operational envelope.

\begin{figure}[ht]
\centering
\includegraphics[width=0.85\textwidth]{blur_vs_altitude}
\caption{Pixel motion blur vs.\ altitude at different flight speeds and exposure times. At the nominal search configuration (\SI{8}{\metre\per\second} at 35\,m altitude, \SI{2}{\milli\second} exposure), blur remains below 1\,px for all altitudes above \SI{15}{\metre}. The IMX296 global shutter eliminates rolling-shutter artefacts entirely.}
\label{fig:blur_vs_altitude}
\end{figure}

\begin{figure}[ht]
\centering
\includegraphics[width=0.75\textwidth]{det_vs_speed}
\caption{Detection probability per flyover vs drone speed at \SI{35}{\metre} altitude, for single-frame detection rates of 0.70, 0.85, and 0.95.  Even at 15\,m/s, the probability exceeds 99\% for $p_d \geq 0.85$.}
\label{fig:det_vs_speed_ext}
\end{figure}

\paragraph{Altitude-dependent speed scheduling.}
To balance coverage rate against detection reliability, \texttt{config.py} implements a linear speed schedule: \SI{6}{\metre\per\second} below \SI{20}{\metre} (maximising frame count during low-altitude centering) and \SI{10}{\metre\per\second} above \SI{50}{\metre} (maximising area coverage during high-altitude search).  This is configured via \texttt{speed\_for\_altitude()} and ensures that the detection frame budget remains above 10 frames per flyover at all operational altitudes.

% ═══════════════════════════════════════════════════════════════
\subsection{Inference Latency Breakdown}\label{app:cv:latency}

Figure~\ref{fig:latency_breakdown_ext} presents the latency breakdown as a stacked bar chart, emphasising that model inference dominates the pipeline at 94\% of total time.  This has two implications for optimisation:

\begin{figure}[ht]
\centering
\includegraphics[width=0.65\textwidth]{latency_breakdown}
\caption{Latency breakdown of the detection pipeline.  Model inference (XNNPACK, 4-thread) accounts for 94\% of total latency; all other stages are negligible in comparison.}
\label{fig:latency_breakdown_ext}
\end{figure}

\begin{enumerate}
    \item \textbf{Optimising non-inference stages yields negligible gains.}  Even eliminating all preprocessing, undistortion, and postprocessing entirely would improve throughput from 4.8 to 5.1\,FPS---a 6\% improvement.  Meaningful speedups require reducing inference time itself, via backend change (NCNN), quantisation (FP16), or hardware acceleration (Hailo-8L NPU).
    \item \textbf{The pipeline is naturally suited to pipelining.}  Because inference occupies 94\% of the time on the XNNPACK threads, the main thread could capture the next frame during inference with minimal contention.  However, as discussed in Section~\ref{sec:cv:inference:pipeline}, this optimisation was deferred due to shared-state complexity in the safety-critical flight loop.
\end{enumerate}
